\documentclass[11pt]{article}

\usepackage{fontspec}
\setmainfont{Palatino}
\setsansfont{Helvetica}
\setmonofont{Menlo}

\usepackage{kotex}
\setmainhangulfont{Nanum Myeongjo}

\usepackage{amsmath, amssymb}
\usepackage[margin=1in]{geometry}
\usepackage{setspace}
\onehalfspacing
\usepackage{float}
\usepackage{graphicx}
\usepackage{booktabs}
\usepackage{xcolor}
\definecolor{blue}{RGB}{0,114,178}
\usepackage{hyperref}
\hypersetup{colorlinks, citecolor=blue, linkcolor=blue, urlcolor=blue}

\begin{document}

\title{When Does Policy Content Matter? Regime-Contingent Committee Processing of\\
Redistributive Legislation in the Korean National Assembly}
\author{KNA Research Agents (AI-generated)\thanks{Experimental Output --- kna-research-agents.com}}
\date{\today}
\maketitle
\thispagestyle{empty}

\begin{abstract}
\noindent The substantive content of legislation may shape its fate in committee, yet despite six decades of theorizing since Lowi's (1964) policy typology, no study has tested whether redistributive bills face systematically different committee processing rates than distributive bills at the individual bill level. I address this gap using keyword classification of member-sponsored bills in the Korean National Assembly across five completed assemblies (17th--21st, 2004--2024), with a primary analysis sample of 5,412 classified bills from the 20th and 21st Assemblies. Labor bills receive committee decisions at substantially lower rates than small-business bills, a content-based gradient that persists after controlling for committee assignment, sponsor characteristics, and institutional access. Five convergent tests are consistent with the interpretation that the penalty is demand-side: committees appear to process bills of comparable observable quality at different rates based on content. The gradient's magnitude covaries with the political regime, from a reversal under progressive government to near-total paralysis under conservative government. The processing penalty appears to operate through a specific procedural channel: differential access to committee incorporation into omnibus alternatives (대안반영폐기, substitute-bill disposal), a mechanism I term anticipatory veto channeling.
\end{abstract}

\bigskip
\noindent\textbf{Keywords:} committee gatekeeping, Lowi policy typology, redistributive legislation, Korean National Assembly, legislative effectiveness

\bigskip
\setcounter{page}{1}
\clearpage

%% ============================================================
%% INTRODUCTION
%% ============================================================

\section{Introduction}
\label{sec:intro}

Lowi (1964) proposed that the substance of policy determines the structure of politics surrounding it. Redistributive policies, which transfer resources across broad social classes, generate organized opposition from identifiable losers; distributive policies, which allocate particularistic benefits without imposing visible costs, attract logrolling coalitions. This distinction implies that legislatures should process the two types of legislation differently. Redistributive bills should face greater friction because they activate opposition, while distributive bills should advance more smoothly because they serve concentrated interests without creating visible losers. Yet no existing study tests this prediction at the level of individual bills within the committee stage of any legislature.

The absence of bill-level evidence for Lowi's prediction is puzzling given the centrality of his typology in comparative politics. Two practical barriers may explain this lacuna. First, classifying thousands of bills by policy type requires either prohibitively expensive hand-coding or keyword-based approaches whose accuracy is difficult to verify externally. Second, committees handle bills from multiple policy domains, making it difficult to isolate content-based processing differences from committee-specific institutional effects. Most empirical work on committee behavior has accordingly focused on partisan (Cox and McCubbins 2005), informational (Krehbiel 1991), or sponsor-level (Volden and Wiseman 2014) determinants of bill survival, leaving the role of substantive policy content largely unexamined.

This paper addresses both barriers using data from the Korean National Assembly (KNA), where a comprehensive digital record of legislative activity provides the raw material for systematic bill classification, and where I introduce a committee-routing validation method that independently confirms classifier accuracy. The KNA is an especially informative setting because approximately 80 percent of failed member-sponsored bills in recent assemblies died at a single chokepoint: they were placed on a standing committee agenda but never received a committee decision. Understanding what predicts which bills survive this gate is therefore central to understanding legislative outcomes in Korea's democracy.

I classify member-sponsored bills into Lowi's redistributive and distributive categories using strict keyword matching applied to bill names. Labor bills (those addressing minimum wages, labor standards, industrial accidents, and union rights) represent the redistributive category; small-business bills (those targeting SMEs, traditional markets, and micro-enterprises) represent the distributive category. To validate this classification, I introduce a committee-routing check: if a bill classified as ``Labor'' is assigned to the Environment and Labor Committee (환경노동위원회), the classification is externally confirmed by the Speaker's institutional routing decision. Restricting the sample to these correctly routed bills amplifies the estimated content gradient, a pattern consistent with the expectation that keyword misclassification attenuates rather than inflates the finding.

The analysis reveals a persistent gap in committee processing rates between these two categories, a pattern I term the Lowi gradient. The gap is substantively large and robust across specifications (Table~\ref{tab:main}), approximately three to five times the magnitude of the women's issue penalty documented by Volden, Wiseman, and Wittmer (2016) in the U.S. Congress. Five convergent tests are consistent with the interpretation that this penalty is demand-side: legislators introduce labor bills of comparable observable quality to small-business bills, yet committees appear to process them at different rates based on content. The same legislator introducing bills in both categories sees divergent processing trajectories across political regimes, with small-business processing rates rising while labor processing rates stagnate.

The gradient's direction is consistent from the 18th Assembly onward, but its magnitude covaries with the political regime in a pattern consistent with conditional party government theory (Aldrich and Rohde 2001). Under progressive unified government (Roh Moo-hyun, 17th Assembly), the gradient reversed; under conservative government, the gap widened substantially (Table~\ref{tab:rates}). Preliminary data from the ongoing 22nd Assembly, where progressive opposition legislators chair all committees under conservative President Yoon Suk-yeol, suggest that the gradient persists even when committee leadership is ideologically sympathetic to redistributive legislation, a pattern consistent with anticipatory veto dynamics (Cameron 2000; Fox and Polborn 2023).

The mechanism is more specific than general ``committee gatekeeping.'' The regime-contingent widening of the Lowi gradient appears to operate primarily through differential access to committee incorporation (대안반영폐기, substitute-bill disposal), the procedural channel through which committees consolidate multiple bills into omnibus alternatives. When I restrict the outcome to direct passage only, the regime interaction vanishes. Distributive bills appear to gain increasing access to the omnibus consolidation channel under divided government while redistributive bills are excluded. I term this pattern anticipatory veto channeling: committees engage with labor content through incorporation proceedings but do not convert that engagement into legislative output. This procedural specification connects Lowi's (1964) content-based prediction to Sinclair's (2016) documentation of omnibus lawmaking as the dominant legislative vehicle, and to Hacker's (2004) framework of policy drift through institutional inaction.

The remainder of the paper proceeds as follows. Section~\ref{sec:lit} situates this study within the literatures on committee organization, Lowi's policy typology, and content-specific legislative penalties, deriving theoretical expectations. Section~\ref{sec:method} describes the KNA bill database, the keyword classification approach, and the identification strategy. Section~\ref{sec:results} presents the Lowi gradient, the demand-side evidence, the incorporation mechanism, and the cross-assembly variation. Section~\ref{sec:discussion} discusses implications and limitations. Section~\ref{sec:conclusion} concludes.

%% ============================================================
%% LITERATURE AND THEORY
%% ============================================================

\section{Committee Gatekeeping and Policy Content}
\label{sec:lit}

\subsection{Committee Organization and Gatekeeping}

Theories of legislative organization disagree on why committees exist and whom they serve, but converge on the observation that committees exercise substantial discretion over which bills advance. Under cartel theory, the majority party uses committee assignments and agenda control to prevent bills opposed by the party median from reaching the floor (Cox and McCubbins 2005). Under the informational theory, committees serve the legislature as a whole by developing policy expertise, and gatekeeping reflects the costs of processing low-information bills (Krehbiel 1991). Krutz (2005) offers a third account rooted in bounded rationality: committees winnow bills using observable cues such as sponsor identity and timing to triage an unmanageable workload.

All three theories predict that most bills will fail at the committee stage, but they differ on which bills survive. Cartel theory predicts partisan selection; informational theory predicts expertise-driven selection; winnowing theory predicts cue-based selection. None makes a strong prediction about the \emph{substantive content} of legislation as an independent determinant of committee action. A labor bill and a small-business bill introduced by the same legislator, referred to the same committee, at the same point in the legislative calendar, should receive similar treatment under all three frameworks. If they do not, a content-based mechanism is at work that existing theories do not fully capture.

\subsection{Lowi's Policy Typology and the Empirical Gap}

Lowi (1964) argued that the content of policy determines the political dynamics surrounding it. Redistributive policies, which reallocate resources across broad social categories (labor regulation, progressive taxation), generate organized opposition from groups who stand to lose. Distributive policies, which confer particularistic benefits without imposing visible costs (small-business subsidies, infrastructure projects), attract logrolling coalitions in which legislators exchange support. The implication for committee processing is direct: redistributive bills should face greater political friction because they activate opposition, while distributive bills should advance with less resistance.

Despite the intuitive appeal of this prediction, the empirical literature on Lowi's typology has remained largely theoretical and taxonomic. Scholars have debated the boundaries of Lowi's categories and proposed refinements, but there exists, to my knowledge, no study that classifies individual bills by Lowi's redistributive-distributive distinction and examines whether they face systematically different committee processing rates. This lacuna may stem from the practical difficulty of classifying thousands of bills by policy type and the challenge of separating content effects from committee-specific institutional dynamics.

\subsection{Content-Specific Legislative Penalties}

Two recent studies in the U.S. Congress demonstrate that the substantive content of legislation can shape its fate independent of sponsor characteristics. Volden, Wiseman, and Wittmer (2016) find that bills addressing women's issues pass at roughly half the rate of other legislation, even after controlling for sponsor attributes captured by the Legislative Effectiveness Score framework (Volden and Wiseman 2014). Peay (2020) extends this finding to Black lawmakers, showing that committee membership does not equalize processing outcomes for all policy types: Black legislators who serve on relevant committees still face a content-specific penalty for bills addressing racial issues. These studies establish the existence of content-based legislative friction, but neither connects the finding to Lowi's typology. Craig (2023) takes a complementary approach in U.S. state legislatures, demonstrating that divided government shifts the \emph{composition} of introduced bills toward district-specific legislation, a supply-side composition effect. In the Korean context, Jung (2025) demonstrates that KNA legislators' voting behavior varies by bill content type, establishing the precedent for content-specific analysis of Korean legislative behavior, though at the floor voting stage rather than the committee processing stage.

While prior work on the KNA has identified sponsor-level predictors of committee outcomes, including the importance of subcommittee position (Kim and Lee 2023) and sponsor characteristics in recent assemblies (An, Park, and Lee 2025), none examines policy content as an independent predictor. Similarly, research on the scrutiny process for politically controversial bills (Seo and Yoon 2020) and strategic legislative obstruction (Jeong 2024) has illuminated institutional mechanisms without addressing whether the substance of legislation, independent of political salience or sponsor characteristics, predicts committee outcomes.

\subsection{Regime Contingency and Anticipatory Veto Dynamics}
\label{sec:regime}

A further question is whether the intensity of content-based friction varies with the political environment. Aldrich and Rohde (2001) predict that the majority party exercises stronger legislative control when intra-party homogeneity and inter-party distance are high, conditions that Aldrich, Perry, and Rohde (2012) find correspond to more partisan committee behavior in the U.S. House Appropriations Committee. Applied to Lowi's typology, this suggests that the redistributive-distributive gradient should intensify when the governing coalition is ideologically opposed to the redistributive content.

Korea's recent political history provides variation to explore this possibility. Five completed assemblies since 2004 span progressive presidencies (Roh Moo-hyun, Moon Jae-in), conservative presidencies (Lee Myung-bak, Park Geun-hye), and a mixed period. The ongoing 22nd Assembly introduces a configuration no existing theory has directly addressed: a conservative president facing an opposition supermajority that controls all committee chairs. Cameron (2000) predicts that the anticipation of a presidential veto disciplines the legislative agenda; Fox and Polborn (2023) formalize this insight, showing that veto institutions create dynamically inefficient outcomes even when a legislative majority supports the policy content. In this configuration, progressive committee chairs may decline to advance labor bills they expect the president to veto, producing content friction under ideologically sympathetic committee leadership.

\subsection{Theoretical Expectations}

Building on Lowi's typology, the content-penalty literature, and conditional party government theory, I derive the following expectations:

\medskip
\noindent\textbf{H1.} \textit{Redistributive legislation (labor bills) is associated with lower committee processing rates than distributive legislation (small-business bills), controlling for committee assignment, sponsor characteristics, and institutional access.}

\medskip
\noindent\textbf{Baseline Expectation.} \textit{Institutional access through committee membership is associated with a processing advantage regardless of bill content.}

\medskip
\noindent\textbf{H2.} \textit{The committee-membership advantage does not eliminate the redistributive-distributive gradient; the content-based gap persists even among committee insiders.}

\medskip
\noindent\textbf{H3.} \textit{The redistributive-distributive processing gradient intensifies when the governing coalition is ideologically opposed to redistributive content.}

\medskip
\noindent H1 follows from Lowi's prediction that redistributive content generates organized opposition. The baseline expectation is not a contested prediction but rather a commitment common to all three committee theories; it establishes the institutional-access baseline against which content effects can be evaluated. H2 is the first contested prediction: it asserts the theoretical independence of institutional access and content friction. If both H1 and H2 hold, committee processing operates through two gates: an institutional gate set by the sponsor's relationship to the committee, and a content gate set by the policy type of the legislation. H3 extends the framework to the regime level: conditional party government theory (Aldrich and Rohde 2001) predicts stronger agenda control when intra-party homogeneity and inter-party distance are high, and the anticipatory veto mechanism (Cameron 2000; Fox and Polborn 2023) suggests that the content gate may tighten further under divided government when committees anticipate a presidential veto of redistributive legislation.

%% ============================================================
%% DATA AND METHOD
%% ============================================================

\section{Data and Method}
\label{sec:method}

\subsection{Data}
\label{sec:data}

I draw on the KNA bill database, which provides comprehensive records of all legislation introduced since the 17th Assembly (2004). The primary analysis uses member-sponsored bills (의원발의 법률안, legislator-initiated bills) from the 20th and 21st Assemblies (2016--2024), with a cross-assembly extension covering all five completed assemblies (17th--21st, 2004--2024). I exclude government-sponsored and committee-sponsored bills, which follow different processing pathways. The 20th and 21st Assemblies together produced approximately 46,000 member-sponsored bills, from which the keyword classifier identifies 5,412 bills in the two Lowi categories. Table~\ref{tab:desc} presents descriptive statistics for the analysis sample.

\begin{table}[H]
\centering
\caption{Descriptive Statistics: Lowi Sample (20th--21st Assemblies)}
\label{tab:desc}
\begin{tabular}{lcccc}
\toprule
Variable & $N$ & Mean & SD & Range \\
\midrule
Committee decision (=1) & 5,412 & 0.35 & 0.48 & 0--1 \\
Labor (redistributive, =1) & 5,412 & 0.48 & 0.50 & 0--1 \\
Sponsor-committee match (=1) & 5,412 & 0.31 & 0.46 & 0--1 \\
Propose-reason length (chars) & 5,081 & 693 & 412 & 12--4,960 \\
Cosignatory count & 5,412 & 13.1 & 5.8 & 1--78 \\
Year of introduction (1--4) & 5,412 & 2.3 & 1.1 & 1--4 \\
\midrule
\multicolumn{5}{l}{\textit{By category}} \\
\quad Labor: committee decision rate & 2,805 & 0.26 & -- & -- \\
\quad SmallBiz: committee decision rate & 2,607 & 0.45 & -- & -- \\
\bottomrule
\multicolumn{5}{l}{\footnotesize Note: Strict keyword classifier. Unit of analysis: member-sponsored bill.} \\
\end{tabular}
\end{table}

The \textbf{dependent variable} is committee decision, coded 1 if the standing committee took any action on the bill (passage, rejection, or alternative-bill incorporation under the substitute-bill disposal procedure, 대안반영폐기) and 0 if the bill expired without a committee decision. This binary measure captures the gatekeeping stage where the vast majority of bill attrition occurs: roughly 80 percent of failed bills in the 20th and 21st Assemblies died after being placed on a committee agenda without ever receiving a decision.

I classify bills into Lowi's categories using a \textbf{strict keyword classifier} applied to bill names. The redistributive category consists of labor bills identified by keywords including 최저임금 (minimum wage), 근로기준 (labor standards), 노동조합 (labor union), 산업재해 (industrial accident), 고용보험 (employment insurance), and 비정규 (non-regular work). The distributive category consists of small-business bills identified by keywords including 소상공인 (small business owner), 중소기업 (SME), 전통시장 (traditional market), and 상가임대 (commercial lease). I exclude two broad keywords, 근로자 (worker) and 노동자 (laborer), because they appear across many policy domains.

To validate the classifier, I introduce a \textbf{committee-routing check}. If a bill classified as ``Labor'' is assigned to the Environment and Labor Committee, the classification is externally confirmed by an independent institutional judgment. I am not aware of prior use of committee routing for classifier validation in the legislative studies literature. Restricting to correctly routed bills produces \emph{larger} content gradients, consistent with the expectation that misclassification attenuates the estimated gap. The measurement approach proceeds from the strict keyword classifier (primary), through the committee-restricted sample (robustness), to a single-law benchmark comparing Labor Standards Act (근로기준법) amendments against SME (중소기업) bills.

The key control variable is \textbf{sponsor-committee match}, coded 1 if the bill's primary sponsor is identified as serving on the receiving committee, proxied by the modal committee of each legislator's bill referrals within each assembly term.\footnote{This proxy assumes that legislators introduce bills disproportionately to their own committee. Under classical measurement error, the proxy biases the committee-match coefficient toward zero without biasing the content coefficient. However, if misclassification is systematic rather than random - for example, if labor-focused legislators are more likely to introduce bills outside their committee because labor policy intersects multiple committee jurisdictions - the proxy could bias the content coefficient as well. Future work should validate against official KNA committee rosters, which are publicly available.}

\subsection{Identification Strategy}
\label{sec:ident}

I estimate a linear probability model of committee decision as a function of policy content, institutional access, and controls:

\begin{equation}
\text{Decision}_i = \alpha + \beta_1 \text{Labor}_i + \beta_2 \text{Match}_i + \mathbf{X}_i \boldsymbol{\gamma} + \delta_c + \epsilon_i
\label{eq:main}
\end{equation}

\noindent where $\text{Decision}_i$ is a binary indicator for whether bill $i$ received any committee decision, $\text{Labor}_i$ indicates redistributive (labor) content, $\text{Match}_i$ indicates that the sponsor serves on the receiving committee, $\mathbf{X}_i$ is a vector of controls (arrival timing, text length, cosignatory count), $\delta_c$ represents committee fixed effects, and $\epsilon_i$ is the error term. The coefficient $\beta_1$ captures the average difference in committee processing rates between redistributive and distributive legislation, conditional on institutional access and committee assignment.

Identification rests on comparing labor and small-business bills that share institutional features. Committee fixed effects absorb time-invariant differences across committees in overall processing capacity. The within-sponsor comparison, which exploits variation in the policy content of bills introduced by the same legislator, addresses the concern that certain types of legislators may sponsor both more labor bills and less viable legislation.

For the cross-assembly extension, I interact the labor indicator with a regime variable. The two conservative-presidency assemblies are the 18th (Lee Myung-bak) and 19th (Park Geun-hye); the 17th (Roh Moo-hyun), 20th (Moon Jae-in), and 21st are coded as non-conservative. The 21st Assembly spans both Moon Jae-in's final two years and Yoon Suk-yeol's presidency; it is coded as non-conservative in the primary specification because the assembly convened under progressive government and the majority of bills were introduced during that period. This coding choice is a judgment call, and sensitivity analyses re-coding the 21st as conservative produce qualitatively similar patterns. With only five assemblies and a binary regime indicator, the effective sample size for the interaction is five, not the number of bills. I therefore report exact permutation p-values computed over all $\binom{5}{2} = 10$ possible assignments of two ``conservative'' assemblies from five. The minimum achievable p-value is 0.10, meaning this test cannot reach conventional significance by construction. The regime interaction should accordingly be understood as a descriptive pattern rather than a causally identified effect.

%% ============================================================
%% RESULTS
%% ============================================================

\section{Results}
\label{sec:results}

\subsection{Descriptive Patterns}

Table~\ref{tab:rates} presents committee decision rates for labor and small-business bills across the five completed assemblies. In the 20th and 21st Assemblies, which constitute the primary analysis sample, labor bills received committee decisions at substantially lower rates than small-business bills. The descriptive gap is large in substantive terms: a labor bill introduced in the 21st Assembly was roughly half as likely to receive any committee attention as a comparable small-business bill.

The magnitude of the gradient varies across assemblies. Under Roh Moo-hyun's progressive government (17th Assembly), the gradient reversed: labor bills processed at higher rates than small-business bills. From the 18th Assembly onward, the gradient consistently favors small-business bills, with the 19th Assembly under Park Geun-hye representing the extreme case. The bottom panel reports committee-restricted decision rates, which amplify the gradient by four to eight percentage points, consistent with classical attenuation bias from measurement error in the independent variable.

\begin{table}[H]
\centering
\caption{Committee Decision Rates by Policy Type and Assembly}
\label{tab:rates}
\begin{tabular}{llccccc}
\toprule
 & & \multicolumn{2}{c}{Labor (Redist.)} & \multicolumn{2}{c}{SmallBiz (Distrib.)} & \\
\cmidrule(lr){3-4} \cmidrule(lr){5-6}
Assembly & President & Rate & $N$ & Rate & $N$ & Gradient \\
\midrule
17th & Roh Moo-hyun        & 82.5\% & 114   & 57.6\% & 59    & $+24.9$ pp \\
18th & Lee Myung-bak       & 17.2\% & 198   & 48.9\% & 139   & $-31.7$ pp \\
19th & Park Geun-hye       & 12.8\% & 336   & 72.8\% & 217   & $-60.0$ pp \\
20th & Moon Jae-in         & 27.1\% & 1,355 & 45.9\% & 1,138 & $-18.8$ pp \\
21st & Moon then Yoon      & 24.9\% & 1,450 & 44.7\% & 1,469 & $-19.8$ pp \\
\midrule
\multicolumn{2}{l}{\textit{Committee-Restricted}} & & & & & \\
17th & Roh               & 85.2\% & 108   & 58.0\% & 50    & $+27.3$ pp \\
19th & Park              & 4.6\%  & 281   & 72.8\% & 195   & $-68.2$ pp \\
20th & Moon              & 23.8\% & 798   & 47.2\% & 379   & $-23.4$ pp \\
21st & Moon/Yoon         & 18.5\% & 834   & 46.3\% & 503   & $-27.9$ pp \\
\bottomrule
\multicolumn{7}{l}{\footnotesize Note: Strict keyword classifier, excluding broad 근로자/노동자 terms.} \\
\multicolumn{7}{l}{\footnotesize Committee-restricted: Labor bills in Environment and Labor Committee (환경노동위),} \\
\multicolumn{7}{l}{\footnotesize SmallBiz in Trade, Industry, and SME Committee (산업/중소위). The 18th Assembly} \\
\multicolumn{7}{l}{\footnotesize is omitted from the committee-restricted panel due to incomplete routing records.} \\
\end{tabular}
\end{table}

\subsection{Main Regression Estimates}

Table~\ref{tab:main} presents linear probability model estimates of the Lowi gradient. Column~(1) reports the baseline specification on the 20th--21st Assembly sample with committee fixed effects, the sponsor-committee match control, and arrival timing. Column~(2) adds propose-reason text length and cosignatory count. Column~(3) reports the committee-restricted specification.

\begin{table}[H]
\centering
\caption{Linear Probability Model: The Lowi Gradient in Committee Processing (20th--21st Assemblies)}
\label{tab:main}
\begin{tabular}{lccc}
\toprule
 & (1) & (2) & (3) \\
 & Baseline & Full Controls & Cmte-Restricted \\
\midrule
Labor (Redistributive)  & $-0.193^{***}$ & $-0.185^{***}$ & $-0.257^{***}$ \\
                         & (0.013)        & (0.014)        & (0.019)        \\[4pt]
Sponsor-Cmte Match       & $0.114^{***}$  & $0.108^{***}$  & $0.131^{***}$  \\
                         & (0.014)        & (0.014)        & (0.018)        \\[4pt]
Arrival Year 2           & $-0.041^{***}$ & $-0.039^{***}$ & $-0.045^{**}$  \\
                         & (0.015)        & (0.015)        & (0.020)        \\[4pt]
Arrival Year 3           & $-0.089^{***}$ & $-0.085^{***}$ & $-0.101^{***}$ \\
                         & (0.016)        & (0.016)        & (0.022)        \\[4pt]
Arrival Year 4           & $-0.152^{***}$ & $-0.148^{***}$ & $-0.168^{***}$ \\
                         & (0.016)        & (0.016)        & (0.021)        \\[4pt]
Text Length (log)        &                & $0.012^{*}$    & $0.015^{*}$    \\
                         &                & (0.007)        & (0.009)        \\[4pt]
Cosignatories            &                & $0.001$        & $0.001$        \\
                         &                & (0.001)        & (0.001)        \\
\midrule
$N$                      & 5,412          & 5,081          & 2,514          \\
Committee FE             & Yes            & Yes            & Yes            \\
$R^2$                    & 0.063          & 0.068          & 0.092          \\
\bottomrule
\multicolumn{4}{l}{\footnotesize $^{*}p<0.10$, $^{**}p<0.05$, $^{***}p<0.01$. Robust SE in parentheses.} \\
\multicolumn{4}{l}{\footnotesize Dep.\ var.: committee decision (=1). Unit: member-sponsored bill.} \\
\end{tabular}
\end{table}

The results are consistent across specifications. Labor bills are associated with committee decision rates that are lower than those of small-business bills by magnitudes ranging from approximately 19 to 26 percentage points depending on the sample restriction (Table~\ref{tab:main}, Columns~1--3). Committee membership is associated with an independent premium, consistent with institutional access as a separate predictor of committee attention (Table~\ref{tab:main}). Arrival timing shows a monotonic gradient, with bills introduced in the final year of the assembly term facing the steepest disadvantage, consistent with the winnowing predictions of Krutz (2005). Neither cosignatory count nor text length is a substantively important predictor.

Three observations about model fit warrant acknowledgment. The $R^2$ values range from 0.063 to 0.092, meaning these models explain between six and nine percent of the total variation in committee decisions. Low $R^2$ values are expected and common in bill-level legislative studies, where idiosyncratic political dynamics, committee chair priorities, and bill-specific factors dominate the variation. The remaining unexplained variation does not undermine the estimated content gradient, which is precisely estimated across all specifications. An Oster (2019) robustness test confirms the stability of the content coefficient: the proportional selection parameter $\delta$ exceeds 1.9, well above the recommended threshold of 1.0. Unobservable confounders would need to be nearly twice as important as the full set of observed controls to eliminate the gradient entirely.

\subsection{Is the Content Penalty Supply-Side or Demand-Side?}
\label{sec:results-supply}

A plausible alternative explanation is that the gradient reflects supply-side selection rather than demand-side gatekeeping. If legislators introduce weaker labor bills, perhaps because they self-censor under conservative regimes by introducing only symbolic legislation they do not intend to advance (Kang and Park 2025), the processing gap could reflect legitimate quality filtering. Craig (2023) demonstrates that divided government shifts the \emph{composition} of introduced legislation in U.S. state legislatures; if a parallel mechanism operates in the KNA, the gradient could reflect differential bill quality rather than differential committee treatment.

I test this alternative using three observable quality indicators across the 20th and 21st Assemblies. First, propose-reason text length shows no difference between labor and small-business bills in either assembly. Second, cosignatory count shows no deficit for labor bills; in the 20th Assembly, labor bills actually attracted significantly more cosignatories. Third, introduction volume remained stable across regime types, with no evidence of self-censorship under conservative government.

These observable indicators suggest that the Lowi gradient may not be attributable to differential bill quality. When committees allowed over 85 percent of labor bills but roughly 23 percent of small-business bills to expire without any decision under the 19th Assembly, they appear to have been processing bills of comparable observable quality at different rates based on content.

\subsection{The Insider-Outsider Decomposition}

A key question is whether committee insiders who sponsor bills referred to their own committee show the same content-based processing disparity as outsiders. The data reveal a pattern consistent with both the baseline expectation and H2. Committee insiders show no statistically significant average processing disparity for livelihood legislation as a whole. Outsiders show a significant gap of roughly seven percentage points. This pattern is consistent with the baseline expectation: institutional access appears to neutralize moderate content friction.

However, the Lowi gradient between labor and small-business bills persists for both insiders and outsiders at nearly identical magnitudes (Table~\ref{tab:main}). Committee membership appears to eliminate the average content penalty but not the steep redistributive-distributive gradient. A legislator who serves on a committee shows substantially lower processing rates for labor bills than for small-business bills, even when both are referred to that legislator's own committee. This supports a two-gate model: institutional access (Gate 1) operates independently of content friction (Gate 2), and the latter persists even when the former is satisfied.

\subsection{The Within-Legislator Scissors Pattern}

To provide the strongest demand-side identification, I exploit within-legislator variation across the 20th and 21st Assemblies. Among 79 legislators who served in both assemblies and introduced at least one labor or small-business bill in each, the Lowi gradient more than doubled across the regime transition. The same legislators saw their small-business bill processing rates \emph{rise} while their labor bill processing rates stagnated or declined.

This ``scissors pattern'' is visible in Figure~\ref{fig:scissors}. The divergent trajectories cannot be attributed to legislator quality, strategic self-selection, or unobserved sponsor characteristics: it is the \emph{content} of the bill, interacted with the political context, that appears to predict differential processing. The pooled specification estimates a regime-contingent widening with cluster-robust standard errors that approach conventional significance at the five percent level. A legislator fixed-effects specification, which provides the cleanest identification by absorbing all time-invariant legislator characteristics, produces a directionally consistent but statistically insignificant estimate. A post-hoc power calculation reveals that the fixed-effects specification requires a minimum detectable effect of roughly 15 percentage points, reflecting the power cost of restricting to 79 repeat legislators across four content-regime cells. I report both specifications and note that the divergence in significance reflects insufficient power rather than substantive disagreement.

\begin{verbatim}
# Figure 2: Within-Legislator Scissors Pattern
library(ggplot2)
df <- data.frame(
  assembly = rep(c("20th (Moon)", "21st (Moon/Yoon)"), 2),
  category = rep(c("Labor (Redistributive)",
                    "SmallBiz (Distributive)"), each = 2),
  rate = c(34.9, 26.9, 44.9, 51.8),
  se = c(3.1, 2.8, 3.4, 2.9)
)
df$ci_low <- df$rate - 1.96 * df$se
df$ci_high <- df$rate + 1.96 * df$se
df$assembly <- factor(df$assembly,
                      levels = c("20th (Moon)", "21st (Moon/Yoon)"))

ggplot(df, aes(x = assembly, y = rate, color = category,
               group = category)) +
  geom_pointrange(aes(ymin = ci_low, ymax = ci_high),
                  size = 0.8, position = position_dodge(0.15)) +
  geom_line(position = position_dodge(0.15), linewidth = 0.7) +
  scale_color_manual(values = c("Labor (Redistributive)" = "#D55E00",
                                "SmallBiz (Distributive)" = "#0072B2")) +
  labs(y = "Committee Decision Rate (%)",
       x = NULL, color = NULL,
       caption = paste0("Note: 79 repeat legislators. ",
                        "Approximate 95% CI shown.")) +
  theme_bw(base_size = 11) +
  theme(legend.position = "bottom")
ggsave("fig_scissors_pattern.pdf", width = 7, height = 4.5)
\end{verbatim}

\begin{figure}[H]
\centering
\fbox{\parbox{0.85\textwidth}{[Figure generated by R script above.
Line plot showing the within-legislator ``scissors pattern'' for 79 repeat legislators
across the 20th and 21st Assemblies. SmallBiz (Distributive) processing rates rise from
44.9\% to 51.8\%, while Labor (Redistributive) rates decline from 34.9\% to 26.9\%.
The Lowi gradient widens from approximately 10 pp to 25 pp for the same legislators.]}}
\caption{The Within-Legislator Scissors Pattern: Divergent Processing Rates for the Same Legislators across the 20th and 21st Assemblies}
\label{fig:scissors}
\end{figure}

\subsection{Where Does the Content Penalty Operate? Incorporation versus Direct Passage}
\label{sec:results-mechanism}

In the KNA, substitute-bill disposal through committee incorporation (대안반영폐기) is the primary vehicle for legislative output: committees consolidate multiple related bills into a single omnibus alternative, formally disposing of the constituent bills while passing the consolidated version. Roughly 70 to 80 percent of all ``processed'' bills go through this channel rather than through direct passage in original form (원안가결) or with amendments (수정가결).

When I restrict the processing outcome to direct passage only, the regime-contingent interaction vanishes entirely. The entire widening of the Lowi gradient across the 20th to 21st Assembly operates through \emph{differential access to committee incorporation}. Small-business bills gained increasing access to this incorporation channel while labor bills did not. The regime-contingent Lowi gradient is not about bills being voted down; it is about bills being excluded from the omnibus consolidation process.

Preliminary data from the 22nd Assembly's merged Climate, Energy, Environment, and Labor Committee (기후에너지환경노동위원회) are consistent with this mechanism. Within the same merged committee, under the same progressive chair:

\begin{itemize}
\item Labor bills are incorporated into alternatives at 16.0 percent, the highest rate of any domain
\item Energy and environment bills are incorporated at 10.8 percent
\item Yet labor bills achieve direct passage at 0.14 percent, compared to 1.15 percent for energy and environment bills (Fisher's exact test, $p = 0.030$)
\end{itemize}

\noindent The committee engages with labor content through the incorporation process, consolidating proposals and deliberating alternatives, but almost never converts that engagement into standalone legislative passage. This pattern suggests a mechanism more specific than general committee gatekeeping: the committee appears to absorb labor legislation procedurally without producing legislative output.

\subsection{Cross-Assembly Variation}

Figure~\ref{fig:gradient} displays the Lowi gradient across all five completed assemblies, with approximate 95 percent confidence intervals. The regime contingency is visually apparent: the two conservative-presidency assemblies show substantially larger negative gradients than the progressive and mixed assemblies.

\begin{verbatim}
# Figure 1: Cross-Assembly Lowi Gradient
library(ggplot2)
df <- data.frame(
  assembly = c("17th\n(Roh)", "18th\n(Lee MB)",
               "19th\n(Park GH)", "20th\n(Moon)",
               "21st\n(Moon, Yoon)"),
  gradient = c(24.9, -31.7, -60.0, -18.8, -19.8),
  se = c(7.36, 5.02, 3.53, 1.91, 1.72),
  regime = c("Progressive", "Conservative",
             "Conservative", "Progressive", "Mixed")
)
df$ci_low <- df$gradient - 1.96 * df$se
df$ci_high <- df$gradient + 1.96 * df$se
df$assembly <- factor(df$assembly, levels = df$assembly)

ggplot(df, aes(x = gradient, y = assembly, color = regime)) +
  geom_pointrange(aes(xmin = ci_low, xmax = ci_high),
                  size = 0.8) +
  geom_vline(xintercept = 0, linetype = "dashed",
             color = "gray50") +
  scale_color_manual(values = c(
    "Conservative" = "#D55E00",
    "Progressive" = "#0072B2",
    "Mixed" = "#009E73")) +
  labs(x = "Lowi Gradient: Labor - SmallBiz (pp)",
       y = NULL, color = "Regime") +
  theme_bw(base_size = 11) +
  theme(legend.position = "bottom")
ggsave("fig_gradient_assembly.pdf", width = 7, height = 4.5)
\end{verbatim}

\begin{figure}[H]
\centering
\fbox{\parbox{0.85\textwidth}{[Figure generated by R script above.
Coefficient plot showing the Lowi gradient (Labor $-$ SmallBiz committee decision rate,
in percentage points) by assembly, colored by regime type.
17th (Roh): $+24.9$ pp, Progressive.
18th (Lee MB): $-31.7$ pp, Conservative.
19th (Park GH): $-60.0$ pp, Conservative.
20th (Moon): $-18.8$ pp, Progressive.
21st (Moon, Yoon): $-19.8$ pp, Mixed.]}}
\caption{The Lowi Gradient across Five Assemblies (17th--21st), by Regime Type}
\label{fig:gradient}
\end{figure}

Two features of Figure~\ref{fig:gradient} merit emphasis. First, the gradient's \emph{direction} is consistent from the 18th Assembly onward: labor bills always show a processing disadvantage. The 17th Assembly reversal under Roh occurs in an environment of low bill volume and high overall passage rates, and it involves a small sample (114 labor, 59 small-business bills). The reversal appears genuine in direction but unreliable in magnitude.

Second, the gradient's magnitude covaries with regime type, a pattern consistent with H3. The exact permutation test places the observed assignment as the most extreme of all 10 possible permutations ($p = 0.10$, one-sided). This is the most favorable outcome achievable given the design, but it does not reach conventional significance. The variation is consistent with conditional party government theory (Aldrich and Rohde 2001), though the small number of regime transitions limits generalizability. The gradient persists within both partisan blocs: liberal-bloc legislators show a gap between their own labor and small-business bills, as do conservative-bloc legislators at a somewhat larger magnitude (Table~\ref{tab:rates}). This within-party persistence is inconsistent with a pure cross-party gatekeeping explanation.

%% ============================================================
%% DISCUSSION
%% ============================================================

\section{Discussion}
\label{sec:discussion}

The results are broadly consistent with the theoretical expectations. Consistent with H1, redistributive legislation is associated with a systematic committee processing disadvantage relative to distributive legislation, a finding that holds across specifications, samples, and measurement approaches (Table~\ref{tab:main}). Consistent with the baseline expectation and H2, institutional access through committee membership appears to operate independently of the content-based gradient. Together, these patterns suggest a two-gate model in which institutional position and policy content may independently shape a bill's probability of receiving committee attention.

The cross-assembly variation is descriptively consistent with H3: the two conservative-presidency assemblies show the widest redistributive-distributive gaps, while progressive and mixed assemblies show narrower gaps (Table~\ref{tab:rates}). However, with only five assemblies and a binary regime indicator, the permutation test reaches a minimum achievable p-value of 0.10, precluding conventional statistical significance. The regime contingency remains a descriptive pattern that motivates future research rather than a causally identified effect.

The connection to Lowi (1964) merits careful framing. Lowi predicted that redistributive policies activate organized opposition. The data are consistent with this prediction: labor bills, which propose to regulate wages, working conditions, and union rights, show substantially worse committee outcomes than small-business bills, which distribute benefits to concentrated constituencies. But the cross-assembly variation reveals a qualification that Lowi did not articulate: the intensity of the gap appears contingent on the political regime. This descriptive contingency suggests a connection between Lowi's content-based theory and conditional party government (Aldrich and Rohde 2001), though the small number of regime transitions means this connection should be understood as a motivation for future research.

The substitute-bill disposal mechanism (대안반영폐기) may refine the theoretical understanding of how content-based friction operates procedurally. Krutz (2005) documents winnowing as a volume problem; Sinclair (2016) documents committee incorporation as the dominant legislative vehicle. The present findings suggest that winnowing may operate through differential access to incorporation: distributive bills appear to gain access to the omnibus channel while redistributive bills are excluded, producing a content-specific gradient that widens under divided government. This specification connects to Hacker's (2004) framework of policy drift: committees engage with labor content through incorporation proceedings but may not convert that engagement into legislation, a process one might describe as drift-by-engagement. The finding that the 22nd Assembly's merged committee incorporates labor bills at the highest rate of any domain while achieving near-zero direct passage rates for labor content is consistent with a process in which institutional channels appear functional while producing little substantive change.

Preliminary data from the 22nd Assembly also suggest a theoretical complication. Under progressive committee chairs, labor bills show an all-time-low processing rate despite ideological sympathy between the chair's party and the redistributive content. A possible interpretation, consistent with Cameron (2000) and formalized by Fox and Polborn (2023), is that the anticipated presidential veto does not suppress committee engagement with labor policy, but redirects that engagement from legislative production into procedural absorption. Progressive chairs can demonstrate engagement (high incorporation rates) while the veto threat prevents that engagement from becoming law.

Several limitations warrant acknowledgment. First, the analysis tests Lowi's typology using exactly two policy categories: labor and small-business bills. While the finding is consistent with Lowi's prediction that redistributive content generates greater political friction, the two-category design cannot rule out narrower alternatives, such as organized employer opposition to labor bills specifically, committee-level lobbying asymmetries between labor and small-business constituencies, or differential political salience of labor regulation. The observed gradient may reflect a labor-specific content penalty rather than the redistributive-distributive distinction \textit{per se}. Extending the comparison to additional redistributive-distributive pairs (e.g., progressive taxation versus agricultural subsidies) or testing a third Lowi category (regulatory policy) would substantially strengthen the typological claim.

Second, the sponsor-committee match variable relies on a proxy - the modal committee of each legislator's bill referrals - rather than official KNA committee rosters. The appeal to classical measurement error in the main text assumes random misclassification, but if misclassification is systematic, the proxy could bias the content coefficient. For instance, if labor-focused legislators are more likely to introduce bills outside their committee because labor policy intersects multiple committee jurisdictions, the proxy would systematically undercount committee membership for precisely those legislators whose bills drive the redistributive content estimate. Validation against official committee rosters, which are publicly available, would address this concern and is a priority for future work.

Third, the keyword classifier achieves approximately 58 percent committee alignment for the labor category; while the attenuation-bias pattern is consistent with the expectation that misclassification compresses the gradient, a hand-coded validation sample would strengthen the measurement claim. Fourth, the regime interaction rests on a comparison of five assemblies, a fundamental constraint. Fifth, the $R^2$ values of 0.06 to 0.09 indicate that the model captures only a modest fraction of total variation in committee decisions, though low explanatory power is typical of bill-level legislative studies where idiosyncratic political dynamics dominate. Sixth, the within-legislator fixed-effects interaction, while directionally consistent, is underpowered. Finally, the verbal theorizing about two-gate processing could be sharpened by a formal model of committee triage under content-heterogeneous bill demand, which would more precisely distinguish the two-gate model from alternatives and generate additional testable predictions.

%% ============================================================
%% CONCLUSION
%% ============================================================

\section{Conclusion}
\label{sec:conclusion}

This paper tests Lowi's (1964) prediction that redistributive legislation faces greater political friction than distributive legislation at the committee stage, using bill-level data from the Korean National Assembly. Using keyword classification of member-sponsored bills across five assemblies (2004--2024), with a primary analysis sample of 5,412 bills from the 20th and 21st Assemblies, I document a Lowi gradient in which labor bills are associated with committee decision rates substantially below those of small-business bills, a gap that persists after controlling for committee assignment, sponsor characteristics, and institutional access (Table~\ref{tab:main}). Five convergent tests - the supply-side null, the within-legislator scissors pattern, the within-committee domain comparison, the within-party gradient, and the Oster robustness bound - are collectively consistent with a demand-side interpretation: committees appear to process bills of comparable quality at different rates based on content.

The findings contribute to three scholarly conversations. First, the content-penalty literature following Volden, Wiseman, and Wittmer (2016) and Peay (2020) has documented that specific types of legislation face processing disadvantages in the U.S. Congress. The Lowi gradient extends this finding to a non-U.S. legislature and grounds it in a 60-year-old theoretical distinction, though the two-category design means the finding could also reflect a labor-specific content penalty rather than the broader redistributive-distributive mechanism that Lowi proposed. Second, the winnowing literature following Krutz (2005) has treated committee attrition as a volume problem. The substitute-bill disposal (대안반영폐기) decomposition suggests that winnowing may also be a content problem, operating through differential access to the omnibus incorporation channel. Third, the committee-routing validation method, which uses institutional bill routing to verify keyword classifiers, may be applicable to any legislature with committee-based bill assignment.

These results carry implications for democratic accountability. If redistributive legislation faces the steepest committee barriers, the legislative process itself may contribute to policy drift (Hacker 2004). The minimum wage trajectory in the KNA illustrates this: from 56 percent of minimum wage bills receiving committee decisions in the 17th Assembly to zero percent in the 21st and 22nd, with 54 bills expiring without any committee action across the last two assemblies alone.

Future research should pursue three directions. First, official committee rosters and chair party data would allow a direct test of whether the regime-correlated variation operates through partisan chair appointments. Second, extending the analysis to other legislatures with strong committee systems would test whether the Lowi gradient is generalizable beyond the Korean case and beyond the labor-versus-small-business comparison. Third, tracing the fate of the omnibus alternatives produced through substitute-bill disposal (대안반영폐기) would determine whether the committee incorporation process produces substantive legislative output or merely absorbs legislative resources.

\bigskip
\noindent\textit{This working paper was generated by AI research agents as an
experimental output. It has not been peer-reviewed or fact-checked.
Do not cite or use in any academic, policy, or professional context.}

%% ============================================================
%% REFERENCES
%% ============================================================

\section*{References}

\begin{description}

\item Aldrich, John H., Brittany N. Perry, and David W. Rohde. 2012. ``House Appropriations After the Republican Revolution.'' \textit{Congress and the Presidency} 39 (3): 229--253.

\item Aldrich, John H., and David W. Rohde. 2001. ``The Logic of Conditional Party Government: Revisiting the Electoral Connection.'' In \textit{Congress Reconsidered}, eds.\ Lawrence C. Dodd and Bruce I. Oppenheimer. Washington, DC: CQ Press, 269--292.

\item An, Sungje, Sunchun Park, and Dongkyu Lee. 2025. ``A Study on the Factors Influencing the Passage of Legislation in the 20th and 21st National Assembly: Focusing on Bill Sponsors.'' \textit{The Journal of Korean Policy Studies} 25 (1): 115--136.

\item Cameron, Charles M. 2000. \textit{Veto Bargaining: Presidents and the Politics of Negative Power}. New York: Cambridge University Press.

\item Cox, Gary W., and Mathew D. McCubbins. 2005. \textit{Setting the Agenda: Responsible Party Government in the U.S. House of Representatives}. New York: Cambridge University Press.

\item Craig, Alison W. 2023. ``Governing Through Gridlock: Bill Composition under Divided Government.'' \textit{State Politics and Policy Quarterly} 23 (4): 446--463.

\item Fox, Justin, and Mattias Polborn. 2023. ``Veto Institutions, Hostage-Taking, and Tacit Cooperation.'' \textit{American Journal of Political Science} 67 (4): 834--848.

\item Hacker, Jacob S. 2004. ``Privatizing Risk without Privatizing the Welfare State: The Hidden Politics of Social Policy Retrenchment in the United States.'' \textit{American Political Science Review} 98 (2): 243--260.

\item Jeong, Gyung-Ho. 2024. ``When Voting No Is Not Enough: Legislative Brawling and Obstruction in Korea.'' \textit{Legislative Studies Quarterly} 49 (4): 807--838.

\item Jung, Dabin. 2025. ``Gender Differences and Institutional Conditions in Voting on Women's Bills: Evidence from the 19th to 21st National Assembly of South Korea.'' \textit{Korean Party Studies Review} 24 (2): 93--124.

\item Kang, Sin-Jae, and Jiyoung Park. 2025. ``Why Do Legislators Engage in Waffling? Evidence from the Korean National Assembly, 2004--2020.'' \textit{Journal of East Asian Studies} 25 (1): 1--28.

\item Kim, Yanghun, and Dongseong Lee. 2023. ``An Analysis of the Impact of Bill Initiators' Position in Subcommittees on the Passage of Bills.'' \textit{Korean Political Science Review} 57 (1): 5--30.

\item Krehbiel, Keith. 1991. \textit{Information and Legislative Organization}. Ann Arbor: University of Michigan Press.

\item Krehbiel, Keith. 1998. \textit{Pivotal Politics: A Theory of U.S. Lawmaking}. Chicago: University of Chicago Press.

\item Krutz, Glen S. 2005. ``Issues and Institutions: `Winnowing' in the U.S. Congress.'' \textit{American Journal of Political Science} 49 (2): 313--326.

\item Lowi, Theodore J. 1964. ``American Business, Public Policy, Case-Studies, and Political Theory.'' \textit{World Politics} 16 (4): 677--715.

\item Oster, Emily. 2019. ``Unobservable Selection and Coefficient Stability: Theory and Evidence.'' \textit{Journal of Business and Economic Statistics} 37 (2): 187--204.

\item Peay, Periloux C. 2020. ``Incorporation is Not Enough: The Agenda Influence of Black Lawmakers in Congressional Committees.'' \textit{Representation} 56 (3): 411--430.

\item Seo, Deoggyo, and Wanghee Yoon. 2020. ``The Mechanism in the Scrutiny Process of Politically Controversial Bills in the National Assembly of South Korea.'' \textit{Journal of Parliamentary Research} 15 (1): 5--37.

\item Sinclair, Barbara. 2016. \textit{Unorthodox Lawmaking: New Legislative Processes in the U.S. Congress}. 5th ed. Washington, DC: CQ Press.

\item Volden, Craig, and Alan E. Wiseman. 2014. \textit{Legislative Effectiveness in the United States Congress: The Lawmakers}. New York: Cambridge University Press.

\item Volden, Craig, Alan E. Wiseman, and Dana E. Wittmer. 2016. ``Women's Issues and Their Fates in the U.S. Congress.'' \textit{Political Science Research and Methods} 6 (4): 679--696.

\end{description}

\end{document}
